\section{2026 Analysis \& Differential Equations}

\subsection{Exam}

\begin{exercise}
\label{analysis:2026:1}
For any $n\in\mathbb{N}$, define
\[
I_n:=\frac{1}{n!}\int_{-\pi/2}^{\pi/2}\left(\frac{\pi^2}{4}-t^2\right)^n\cos t\,dt.
\]
\begin{itemize}
\item[(a)] Prove
\[
I_{n+1}=2(2n+1)I_n-\pi^2I_{n-1}.
\]
\item[(b)] Show that $\pi^2\notin\mathbb{Q}$, that is, $\pi^2$ is irrational.
\end{itemize}
\end{exercise}

\begin{solution}
Put $a=\pi/2$ and
\[
P_n(t):=\frac{(a^2-t^2)^n}{n!}.
\]
Then $I_n=\int_{-a}^a P_n(t)\cos t\,dt$.  A direct differentiation gives, for $n\geq 1$,
\[
P_{n+1}''(t)
=\pi^2P_{n-1}(t)-2(2n+1)P_n(t).
\]
Indeed,
\[
\begin{aligned}
P_{n+1}''(t)
&=-2P_n(t)+4t^2\frac{(a^2-t^2)^{n-1}}{(n-1)!}  \\
&=-2P_n(t)+4a^2P_{n-1}(t)-4nP_n(t) \\
&=\pi^2P_{n-1}(t)-2(2n+1)P_n(t).
\end{aligned}
\]
Since $P_{n+1}(\pm a)=P_{n+1}'(\pm a)=0$, integrating by parts twice yields
\[
\int_{-a}^a P_{n+1}''(t)\cos t\,dt
=-\int_{-a}^a P_{n+1}(t)\cos t\,dt
=-I_{n+1}.
\]
Integrating the displayed identity for $P_{n+1}''$ against $\cos t$ therefore gives
\[
-I_{n+1}=\pi^2I_{n-1}-2(2n+1)I_n,
\]
which is the desired recurrence.

For irrationality, first note that $I_n>0$ for all $n$, because
$a^2-t^2\geq 0$ and $\cos t\geq 0$ on $[-a,a]$, and the integrand is not identically zero.  Also
\[
I_0=\int_{-\pi/2}^{\pi/2}\cos t\,dt=2,
\]
and
\[
\begin{aligned}
I_1
&=\int_{-\pi/2}^{\pi/2}\left(\frac{\pi^2}{4}-t^2\right)\cos t\,dt \\
&=\frac{\pi^2}{4}\cdot 2-\left[t^2\sin t+2t\cos t-2\sin t\right]_{-\pi/2}^{\pi/2}
=4.
\end{aligned}
\]
Suppose, for contradiction, that $\pi^2=p/q$ with positive integers $p,q$.  Define
\[
J_n:=q^n I_n.
\]
Then $J_0=2$ and $J_1=4q$ are integers.  Multiplying the recurrence by $q^{n+1}$ gives
\[
J_{n+1}=2(2n+1)qJ_n-pqJ_{n-1},
\]
so induction gives $J_n\in\mathbb{Z}$ for all $n$.

On the other hand,
\[
0<I_n
\leq \frac{1}{n!}\int_{-\pi/2}^{\pi/2}\left(\frac{\pi^2}{4}\right)^n\,dt
=\frac{\pi}{n!}\left(\frac{\pi^2}{4}\right)^n.
\]
Hence
\[
0<J_n
\leq \frac{\pi}{n!}\left(\frac{q\pi^2}{4}\right)^n
=\frac{\pi}{n!}\left(\frac{p}{4}\right)^n
\longrightarrow 0.
\]
For all sufficiently large $n$, this gives an integer $J_n$ satisfying $0<J_n<1$, a contradiction.  Therefore $\pi^2$ is irrational.
\end{solution}

\begin{exercise}
\label{analysis:2026:2}
For any $n\times n$ complex matrix $X$, define
\[
e^X=\sum_{m\geq 0}\frac{X^m}{m!}.
\]
\begin{itemize}
\item[(a)] If $X,Y$ are $n\times n$ complex matrices satisfying $[X,Y]=0$, then
\[
e^Xe^Y=e^{X+Y}=e^Ye^X.
\]
\item[(b)] If $X,Y$ are $n\times n$ complex matrices satisfying
\[
[X,[X,Y]]=[Y,[X,Y]]=0,
\]
then
\[
e^Xe^Y=e^{X+Y+\frac{1}{2}[X,Y]}.
\]
\end{itemize}
\end{exercise}

\begin{solution}
For part (a), since $X$ and $Y$ commute, every power of $X$ commutes with every power of $Y$.  Therefore the Cauchy product of the two absolutely convergent matrix series is legitimate:
\[
\begin{aligned}
e^Xe^Y
&=\left(\sum_{j\geq 0}\frac{X^j}{j!}\right)
  \left(\sum_{k\geq 0}\frac{Y^k}{k!}\right) \\
&=\sum_{m\geq 0}\sum_{j+k=m}\frac{X^jY^k}{j!k!}
=\sum_{m\geq 0}\frac{1}{m!}\sum_{j=0}^m\binom{m}{j}X^jY^{m-j} \\
&=\sum_{m\geq 0}\frac{(X+Y)^m}{m!}
=e^{X+Y}.
\end{aligned}
\]
The same argument gives $e^Ye^X=e^{Y+X}=e^{X+Y}$.

For part (b), set
\[
C:=[X,Y].
\]
The assumptions say $[X,C]=[Y,C]=0$.  In particular, $C$ commutes with $X$, $Y$, and $X+Y$.

Consider
\[
\Phi(t):=e^{tX}e^{tY}.
\]
We claim that
\[
\Phi'(t)=(X+Y+tC)\Phi(t).
\]
Indeed,
\[
\Phi'(t)=Xe^{tX}e^{tY}+e^{tX}Ye^{tY}.
\]
Since $[X,C]=0$, the conjugation identity truncates after the first commutator:
\[
e^{tX}Ye^{-tX}=Y+t[X,Y]=Y+tC.
\]
Thus
\[
e^{tX}Ye^{tY}=(Y+tC)e^{tX}e^{tY},
\]
and the claimed differential equation follows.

Now define
\[
\Psi(t):=e^{t(X+Y)+\frac{t^2}{2}C}.
\]
Because $C$ commutes with $X+Y$,
\[
\Psi'(t)=(X+Y+tC)\Psi(t).
\]
Also $\Phi(0)=\Psi(0)=I$.  Uniqueness for linear ODEs in the finite-dimensional space of complex matrices gives $\Phi(t)=\Psi(t)$ for all $t$.  Setting $t=1$ yields
\[
e^Xe^Y=e^{X+Y+\frac12[X,Y]}.
\]
\end{solution}

\begin{exercise}
\label{analysis:2026:3}
Let $f(z)$ be a smooth function on $|z|<R$, and fix a real number $\rho$ such that $0<\rho<R$.  Put
\[
g(z)=\frac{1}{2\pi\sqrt{-1}}\iint_{|\zeta|\leq \rho}
\frac{f(\zeta)}{\zeta-z}\,d\zeta\wedge d\bar{\zeta}.
\]
Then
\[
\frac{\partial}{\partial\bar z}g(z)=f(z),\qquad |z|<\rho.
\]
\end{exercise}

\begin{solution}
We use the standard notation
\[
\frac{\partial}{\partial\bar z}
=\frac12\left(\frac{\partial}{\partial x}+i\frac{\partial}{\partial y}\right),
\qquad z=x+iy.
\]
Since $d\zeta\wedge d\bar{\zeta}=-2i\,dA(\zeta)$, where $dA$ is planar Lebesgue measure, the given function can be rewritten as
\[
g(z)=\frac{1}{\pi}\iint_{|\zeta|\leq \rho}
\frac{f(\zeta)}{z-\zeta}\,dA(\zeta).
\]
Thus $g$ is the convolution of $f\mathbf{1}_{|\zeta|\leq\rho}$ with the Cauchy kernel
\[
E(z)=\frac{1}{\pi z}.
\]

The key distributional identity is
\[
\frac{\partial E}{\partial\bar z}=\delta_0.
\]
For completeness, recall why the normalization is correct.  For a test function $\varphi\in C_c^\infty(\mathbb{C})$, applying Green's formula on $\{|z|>\varepsilon\}$ and then using the Cauchy integral formula gives
\[
-\int_{|z|>\varepsilon}\frac{1}{\pi z}
\frac{\partial\varphi}{\partial\bar z}\,dA(z)
=\frac{1}{2\pi i}\int_{|z|=\varepsilon}\frac{\varphi(z)}{z}\,dz+o(1)
\longrightarrow \varphi(0).
\]
Hence $\partial_{\bar z}E=\delta_0$.

Now fix a point $z$ with $|z|<\rho$.  Since $z$ lies in the interior of the disk of integration, the distributional computation is local near $z$:
\[
\begin{aligned}
\frac{\partial g}{\partial\bar z}(z)
&=\iint_{|\zeta|\leq\rho}
f(\zeta)\frac{\partial}{\partial\bar z}
\left(\frac{1}{\pi(z-\zeta)}\right)dA(\zeta) \\
&=\iint_{|\zeta|\leq\rho}f(\zeta)\delta_z(\zeta)\,dA(\zeta)
=f(z).
\end{aligned}
\]
This proves the asserted Cauchy--Green formula on $|z|<\rho$.
\end{solution}

\begin{exercise}
\label{analysis:2026:4}
Let $C([0,1])$ be the space of continuous functions over $[0,1]$.  Consider the $\infty$-norm on $C([0,1])$
\[
\|f\|_{L^\infty}:=\max_{x\in[0,1]}|f(x)|.
\]
Then $(C([0,1]),\|\cdot\|_{L^\infty})$ is a Banach space.  Consider the closed linear subspace $\mathcal{P}$ of $C([0,1])$ consisting of polynomials over $[0,1]$.  Show that $\dim\mathcal{P}<\infty$.
\end{exercise}

\begin{solution}
For each $m\geq 0$, let
\[
\mathcal{P}_m:=\{p\in\mathcal{P}: p \text{ is represented on } [0,1]
\text{ by a polynomial of degree at most } m\}.
\]
Then $\mathcal{P}_m$ is a finite-dimensional linear subspace of $C([0,1])$, hence it is closed in $C([0,1])$.  Therefore $\mathcal{P}_m$ is also closed in the Banach space $\mathcal{P}$.

Because every element of $\mathcal{P}$ is a polynomial, we have
\[
\mathcal{P}=\bigcup_{m=0}^{\infty}\mathcal{P}_m.
\]
The space $\mathcal{P}$ is complete, since it is closed in the Banach space $C([0,1])$.  By the Baire category theorem, a complete metric space cannot be a countable union of closed subsets with empty interior.  Hence some $\mathcal{P}_N$ has nonempty interior in $\mathcal{P}$.

But $\mathcal{P}_N$ is a linear subspace of $\mathcal{P}$.  A proper linear subspace of a normed vector space has empty interior: if it contained a ball around some point, translating would give a ball around $0$, and scaling that ball would force the whole space to lie in the subspace.  Therefore $\mathcal{P}_N=\mathcal{P}$.

It follows that
\[
\dim\mathcal{P}=\dim\mathcal{P}_N\leq N+1<\infty.
\]
\end{solution}

\begin{exercise}
\label{analysis:2026:5}
Consider the heat equation with periodic boundary conditions on the unit interval $[0,1]$
\[
\begin{cases}
\partial_tu(t,x)-\partial_{xx}u(t,x)=0, & \forall t>0,\ \forall x\in(0,1),\\
u(t,0)=u(t,1), & \forall t>0,\\
\partial_xu(t,0)=\partial_xu(t,1), & \forall t>0.
\end{cases}
\tag{1}
\]
\begin{itemize}
\item[(1)] Write an orthonormal basis for $L^2([0,1])$, you do not need to verify this.  Show that the double derivative operator $f\mapsto f''$ with domain
\[
F:=\{f\in H^2([0,1]): f(0)=f(1),\ f'(0)=f'(1)\}\subset L^2([0,1])
\]
is a self-adjoint operator.
\item[(2)] Recall that the fundamental solution of the heat equation (1) refers to a function $h(t,x,y)$ on $(0,\infty)\times[0,1]\times[0,1]$ which satisfies that
\begin{itemize}
\item[(a)] for any $y\in[0,1]$, $(t,x)\mapsto h(t,x,y)$ is a solution of (1);
\item[(b)] for any $x\in[0,1]$, $h(t,x,y)\to\delta_x$ weakly as $t\to 0$.  That is, for any continuous function $\varphi$ on $[0,1]$ with $\varphi(0)=\varphi(1)$,
\[
\int_0^1h(t,x,y)\varphi(y)\,dy\to\varphi(x)
\]
as $t\to 0$.
\end{itemize}
Find the fundamental solution $h(t,x,y)$ of the heat equation (1) and check that it satisfies the conditions (a) and (b).
\item[(3)] For each fixed $t>0$, where does the fundamental solution $h(t,x,y)$ attain its maximum (that is, for what $x,y\in[0,1]$ is $h$ maximum)?  What can you say about $h(t,x,y)$ when $t\to\infty$?  Finally, show that $h(t,x,y)\sim t^{-1/2}\exp(-d(x,y)^2/t)$ for small $t>0$, where
\[
d(x,y):=\min\{|x-y|,1-|x-y|\}.
\]
That is, there exist constants $A,B,C,D>0$, such that for all $t\in(0,1)$, $x,y\in[0,1]$,
\[
\frac{A}{\sqrt{t}}e^{-\frac{Bd(x,y)^2}{t}}
\leq h(t,x,y)\leq
\frac{C}{\sqrt{t}}e^{-\frac{Dd(x,y)^2}{t}}.
\]
\end{itemize}
\end{exercise}

\begin{solution}
\textbf{Part (1).}
Use the complex Fourier basis
\[
e_k(x):=e^{2\pi ikx},\qquad k\in\mathbb{Z}.
\]
This is an orthonormal basis of $L^2([0,1])$.  For $e_k$,
\[
e_k''=-(2\pi k)^2e_k.
\]
If $f=\sum_{k\in\mathbb{Z}}\hat f_k e_k$, then the periodic Sobolev domain $F$ is exactly
\[
F=\left\{f\in L^2([0,1]):
\sum_{k\in\mathbb{Z}}(1+k^4)|\hat f_k|^2<\infty\right\},
\]
and
\[
f''=\sum_{k\in\mathbb{Z}}-(2\pi k)^2\hat f_k e_k.
\]
Under the Fourier transform, the operator $A f=f''$ with domain $F$ becomes multiplication on $\ell^2(\mathbb{Z})$ by the real sequence $-(2\pi k)^2$ on its maximal natural domain
\[
\left\{(\hat f_k)\in\ell^2:
\sum_{k\in\mathbb{Z}}k^4|\hat f_k|^2<\infty\right\}.
\]
A real multiplication operator on this maximal domain is self-adjoint.  Therefore $f\mapsto f''$ with the periodic $H^2$ domain $F$ is self-adjoint.  Equivalently, the integration-by-parts boundary form
\[
\langle f'',g\rangle-\langle f,g''\rangle
=\left[f'\bar g-f\bar g'\right]_{0}^{1}
\]
vanishes for $f,g\in F$, and the Fourier description shows that the symmetric operator has no larger adjoint domain.

\textbf{Part (2).}
The heat kernel is obtained by evolving each Fourier mode:
\[
h(t,x,y)
=\sum_{k\in\mathbb{Z}}e^{-4\pi^2k^2t}e^{2\pi ik(x-y)}.
\]
Equivalently,
\[
h(t,x,y)
=1+2\sum_{k=1}^{\infty}e^{-4\pi^2k^2t}\cos(2\pi k(x-y)).
\]
For every $t>0$ this series is absolutely and uniformly convergent, and the differentiated series are also uniformly convergent on compact subsets of $(0,\infty)\times[0,1]\times[0,1]$.  Thus we may differentiate term by term:
\[
\partial_t h
=\sum_{k\in\mathbb{Z}}(-4\pi^2k^2)e^{-4\pi^2k^2t}e^{2\pi ik(x-y)}
=\partial_{xx}h.
\]
Also each mode is $1$-periodic in $x$, so
\[
h(t,0,y)=h(t,1,y),
\qquad
\partial_xh(t,0,y)=\partial_xh(t,1,y).
\]
This checks condition (a).

The same kernel has the periodized Gaussian representation
\[
h(t,x,y)=\sum_{m\in\mathbb{Z}}\frac{1}{\sqrt{4\pi t}}
\exp\left(-\frac{(x-y+m)^2}{4t}\right),
\]
which follows from the Poisson summation formula.  Let $\varphi$ be continuous on $[0,1]$ with $\varphi(0)=\varphi(1)$, and extend it periodically to a continuous function $\tilde\varphi$ on $\mathbb{R}$.  Then
\[
\int_0^1h(t,x,y)\varphi(y)\,dy
=\int_{\mathbb{R}}\frac{1}{\sqrt{4\pi t}}
\exp\left(-\frac{(x-\eta)^2}{4t}\right)\tilde\varphi(\eta)\,d\eta.
\]
The Gaussian kernels on the real line form an approximate identity, so the last expression tends to $\tilde\varphi(x)=\varphi(x)$ as $t\to 0$.  This proves condition (b).

\textbf{Part (3).}
The kernel depends only on the periodic difference $x-y$.  Write
\[
H_t(r):=\sum_{k\in\mathbb{Z}}e^{-4\pi^2k^2t}e^{2\pi ikr}.
\]
Then $h(t,x,y)=H_t(x-y)$ and
\[
H_t(0)-H_t(r)
=2\sum_{k=1}^{\infty}e^{-4\pi^2k^2t}\bigl(1-\cos(2\pi kr)\bigr)\geq 0.
\]
Equality holds precisely when $r\in\mathbb{Z}$.  Hence, for $x,y\in[0,1]$, the maximum is attained exactly when the two points are the same point on the circle, namely $d(x,y)=0$; in interval coordinates this means $x=y$, together with the endpoint identifications $(x,y)=(0,1)$ and $(1,0)$.

As $t\to\infty$, all nonzero Fourier modes vanish exponentially fast:
\[
h(t,x,y)\longrightarrow 1
\]
uniformly in $x,y$.  This is the convergence to the constant equilibrium with total mass one.

It remains to prove the small-time Gaussian bounds.  Let
\[
r=d(x,y)\in[0,1/2].
\]
Using the periodized Gaussian formula and choosing an integer $m_0$ for which $|x-y+m_0|=r$, we get the lower bound
\[
h(t,x,y)\geq \frac{1}{\sqrt{4\pi t}}\exp\left(-\frac{r^2}{4t}\right).
\]
Thus the desired lower estimate holds, for example, with
\[
A=\frac{1}{\sqrt{4\pi}},
\qquad
B=\frac14.
\]

For the upper bound, write $s=x-y$.  Since $r=\min_{m\in\mathbb{Z}}|s+m|$, every summand satisfies $|s+m|\geq r$.  Therefore
\[
\exp\left(-\frac{(s+m)^2}{4t}\right)
\leq
\exp\left(-\frac{r^2}{8t}\right)
\exp\left(-\frac{(s+m)^2}{8t}\right).
\]
For $0<t<1$,
\[
\sum_{m\in\mathbb{Z}}\exp\left(-\frac{(s+m)^2}{8t}\right)
\leq C_0
\]
with a universal constant $C_0$, because the left side is uniformly bounded for $s\in[0,1]$ and $0<t<1$.  Hence
\[
h(t,x,y)
\leq \frac{C_0}{\sqrt{4\pi t}}
\exp\left(-\frac{r^2}{8t}\right).
\]
This is the desired upper estimate, for instance with
\[
C=\frac{C_0}{\sqrt{4\pi}},
\qquad
D=\frac18.
\]
Combining the two estimates proves
\[
h(t,x,y)\asymp t^{-1/2}\exp\left(-\frac{d(x,y)^2}{t}\right)
\]
in the stated two-sided sense for $0<t<1$.
\end{solution}
